\documentclass[10pt]{article}
\usepackage{lingmacros}
\usepackage{tree-dvips}
\usepackage{makecell}
\begin{document}
\pagenumbering{gobble}
\section*{Enzimi}
I coenzimi derivano per la maggior parte dalle vitamine. Una carenza di vitamina B12 (origine animale) provoca anemia, ma può cronicizzare colpendo anche il SN.
\paragraph*{Inibitori competitivi}\underline{Vmax uguale} se la \underline{[S] aumenta},  quindi \underline{aumenta la Km}. \\Esempi: la succinato deidrogenasi viene inibita dall'ossalacetato. La xantina ossidasi trasforma la xantina in acido urico e nel caso di produzione eccessiva causa \underline{gotta}, viene inibito da allopurinolo. Viene sfruttato fluorouracile come inibitore competitivo nella produzione di DNA delle \underline{cellule tumorali}. 
\paragraph*{Inibitori non competitivi} Formano EI+S e/o ESI, \underline{Km uguale} perché è un sito diverso, \underline{Vmax diminuisce}.
\paragraph*{Inibitori incompetitivi} Formano solo ESI, \underline{Km e Vmax diminuiscono}.
\paragraph*{Allosterici}Effettori omotropici (il substrato stesso, sigmoide) e eterotropici (inibiscono la reazione a monte).
\section*{Bioenergetica}
Una reazione endoergonica non spontanea non diventa spontanea con un catalizzatore. $\Delta$G= G prodotti - G reagenti. I prodotti devono essere più stabili e quindi $\Delta$G negativo. Nelle reazioni a catena si sommano.\\
$\Delta G^0$ è in condizioni "standard" ma non avviene mai nelle cellule perché presuppone un pH 0, la $\Delta G^0_1$ considera il pH 7/7,4 e viene quindi usato per calcolare il $\Delta$G secondo la formula: \textbf{$\Delta$G= $\Delta G^0_1$+2,303RTlog$\frac{[B]}{[A]}$}.\\
La [Mg$^{2+}$] incide sul rilascio di energia dall'idrolisi di ATP, più magnesio è presente meno energia viene liberata.
\subsection*{Redox}
In questo caso per valutare la spontaneità della reazione si una il potenziale di riduzione, misurato attraverso la pila. $\Delta$E= E prodotti - E reagenti; e grazie alla formula negativa ricaviamo $\Delta$G$<0$=$\Delta$E$>0$.\\
Le deidrogenasi NAD dipendenti hanno legami molto più tenaci con il coenzima, rispetto a quelle FAD dipendenti, grazie al Rossman fold una sequenza glicinica che forma ripiegamenti bab.\\
Mentre il NAD$^+$ si riduce a NADH in un unico passaggio, il FAD passa a FADH$_2$ in due passaggi distinti.
\section*{Glicolisi}
Il glucosio viene usato per la produzione di polimeri strutturali, ribosio-5-fosfato (via dei pentosofosfati), glicogeno e piruvato.\\
Nella 4° tappa la formazione di gliceraldeide3P non è spontanea e quindi la 6° tappa avviene istantaneamente abbassando quindi la concentrazione dei prodotti\\
La piruvato chinasi non viene disattivata nel muscolo scheletrico.\\
Non vengono usati i lipidi perché richiedono più tempo, sono ossigeno dipendenti e non possono formare gluconeogenesi. 
\section*{Gluconeogenesi} 
La glucosio-6-fosfatasi si trova nel RE liscio degli epatocini, non essendo presente in muscolo e cervello avviene solo nel fegato. 
\end{document}
